\chapter{An Illustrated Ear Guide to Iconic Instruments}
\label{chap:iconic-instruments}

\classquote{凡音之起，由人心生也。人心之动，物使之然也。}{《史记·乐书》 / Records of the Grand Historian, "Book of Music" (Yue Shu)}{Translation: Sound arises from the human heart, and the heart is moved by things. Note: the "flavor" of a legendary instrument lives both in the machine and in the recording, arrangement, and listener's experience.}

This is not a collector's catalogue, and a brand name is not a promise of
sound quality. The chapter answers three questions only: where an instrument
came from, what makes its sound recognisable, and where it can land in Yeban.
The fourteen instruments are drawn from three histories -- drum machines,
synthesizers, and electric instruments. They are listening references, not a
list of Yeban clones and certainly not a shopping list that promises an
instant re-creation.

Historical notes give the maker, the decade, and publicly verifiable works.
Works are named by title and artist only; no lyrics are quoted. Yeban's
product boundaries follow the live code: there is an SFZ sampler path,
generated drum styles, and six factory presets on the poly synth. There is no
licensed original sample set, no 303 glide knob, no 707 kit, no Linn kit, no
chorus, and no rotary speaker. Where Yeban can approximate, the text says
``approximate''; where it cannot, it says ``missing part''. ``Sounds like''
is never written as ``is''.

\section{Four dimensions for listening}
\label{sec:iconic-listening}

Do not start with a year. Listen through the four dimensions below, then go
back to the brand and model. What you hear is a sound, not a story.

\begin{longtable}{p{0.16\linewidth}p{0.38\linewidth}p{0.36\linewidth}}
\hline
Dimension & What to listen for & Typical example \\
\hline
\endhead
Envelope and transient & How fast the attack is, how long the decay lasts,
and whether the tail crosses the next beat & the long 808 kick, a clean
short Linn hit, the bell-like tail of a Rhodes \\
Spectrum and resonance & Where the fundamental sits, whether the harmonics
are regular or emphasised, and whether the filter ``sings'' & the high
resonance of a 303, the glass of a DX7, the drawbar harmonics of a Hammond \\
Time and motion & Glide, swing, rotation, tape wow & a 303 slide, 909 shuffle,
Leslie rotation, Mellotron tape motion \\
Arrangement role & Whom the instrument replaces in the band: bass, backbeat,
pad, lead, solo & an 808 as bass, a 909 as the four-on-the-floor frame,
an MPC as the sampling and sequencing centre \\
\hline
\end{longtable}

Arrangement role decides ``does it sound right'' before anything else. Put a
long 808 kick into a pop section that needs a clean backbeat and the role is
wrong even if the timbre is perfect. Ask a Hammond to comp like a short,
percussive piano and it will also sound wrong. Set the role first, then the
envelope and spectrum, and only then the motion.

The Yeban landing column in this chapter uses three grades:
\textbf{directly usable} (a factory style or preset already carries the main
traits), \textbf{approximate} (existing devices cover part of the character
and the missing part is named), and \textbf{own material required} (Yeban
does not ship a licensed sample, so the user must supply an SFZ or a
third-party instrument). The grades are never mixed silently.

\section{Drum machines and samplers}
\label{sec:iconic-drum-machines}

The drum machine began as a replacement drummer and became a composer of
electronic drums. The five models below represent three routes: pure analog
(808), analog plus digital samples (909), all samples (707 and Linn), and
sampler plus sequencer (MPC). Historical status does not mean Yeban has a
matching kit; each entry names both the landing point and the gap.

\subsection{TR-808: a failed drum machine that became a bass}
\label{sec:iconic-808}

Roland introduced the TR-808 in 1980. It generated drum sounds with analog
circuits rather than playing samples. It was meant to replace a real drummer,
but sounded too electronic and sold poorly; production ended in 1983 after
about 12,000 units. The second-hand market later carried it into hip-hop,
electro, and early house. ``Planet Rock'' (Afrika Bambaataa, 1982) and
``Sexual Healing'' (Marvin Gaye, 1982) are early examples of its use in
popular recordings.

\textbf{Ear fingerprint}: the kick is not a short ``thump'' but a long,
sine-like decay with a falling pitch; the tail can cross into the next beat.
The snare is noise-heavy; the clap and snare do not chase acoustic realism.
Congas and cowbell carry an audible analog-circuit character. The 808 kick
works as a bass not because it is merely ``low'', but because its
fundamental sits in the bass range and its decay is long enough that
repetition turns it into a bassline.

\textbf{Yeban landing}: directly usable. Create a MIDI track, press F9 for
the device rack, and generate a drum style from \texttt{Trap} (default
140 BPM) or \texttt{Hip-Hop} (default 92 BPM). The style installs an
808-oriented synthetic kit. Keep the kick on 1/3 and the snare on 2/4 in the
step grid, then write the sub line on the \texttt{bass808} lane. For a longer
kick tail, change \texttt{kick\_decay\_s} through the MCP
\texttt{DeviceParam}; the device card has no per-drum knob, so do not look
for an ``808 kick decay'' control in the GUI.

\textbf{Boundary}: \texttt{bass808} is a separate voice; the kick does not
automatically become a bass. The kick supplies a low-frequency pulse and
\texttt{bass808} supplies a sustained bass line. When both are on, choose one
as the low-frequency carrier and make the other yield, or the low end will
fight itself.

\subsection{TR-909: analog body, digital metal}
\label{sec:iconic-909}

Roland introduced the TR-909 in 1983 as the successor to the 808. It kept
analog circuits for the kick and toms, but used digital samples for the hi-hats
and cymbals. It was also Roland's first drum machine with MIDI, and its
sequencer offered flam and shuffle. About 10,000 units were made; it was not a
commercial success either, but later became the backbone of house and techno.
``Love Can’t Turn Around (1986), ``Strings of Life (1987), and
``The Whistle Song (1991) are common references.

\textbf{Ear fingerprint}: the kick is shorter and harder than the 808; the
snare has a clear crack; the hats and cymbals are sampled metal, cleaner and
more clipped than the 808's analog metal. The most quoted rhythmic outline is
a four-on-the-floor kick with an offbeat open hat. The 909 shuffle belongs to
its sequencer, not to the timbre itself.

\textbf{Yeban landing}: approximate, with a usable frame. Choose
\texttt{House} (default 124 BPM), \texttt{Techno} (default 128 BPM), or
\texttt{EDM} (default 128 BPM) in the drum style generator. The generator
installs a 909 base or a 909 variant. In the step grid, make ``four-on-the-
floor kick + snare or clap on 2/4 + offbeat open hat''. House wants a longer
open hat, techno a shorter, harder kick; use the MCP \texttt{DeviceParam} to
adjust \texttt{hihat\_open\_decay\_s} and \texttt{kick\_decay\_s} in those
directions. The GUI has no separate ``909 kit'' selector; style generation is
the landing point.

\textbf{Boundary}: the 909 shuffle is not exposed as a global knob. The
style-table swing and note micro-timing can approximate it, but do not claim
to reproduce the 909 sequencer's feel.

\subsection{TR-707: eight-bit, dry offbeats}
\label{sec:iconic-707}

Roland introduced the TR-707 in 1985. It was Roland's first drum machine to
use sampled sounds throughout, mostly at 8-bit resolution, with some metal
and envelope details at 6-bit. It had no per-voice pitch control, but offered
four global shuffle settings, an LCD, and MIDI/DIN sync. It was common in
early house, especially acid house; ``Acid Tracks'' (Phuture, 1987) and
``Can You Feel It?'' (Mr. Fingers, 1986) both used it.

\textbf{Ear fingerprint}: the closed hat is short, dry, and slightly
quantised; the snare and rimshot are not thick, but the offbeat is crisp.
There is none of the 808's long low end and none of the 909's big cymbal
tail. It is good at cutting the rhythm clean, not at carrying the low end
alone.

\textbf{Yeban landing}: own material or approximation. Yeban has no 707 kit.
The closest factory route is the 909 kit generated by \texttt{House} or
\texttt{Techno}, using its shorter closed hats and dry snare, then shortening
the open-hat decay. To keep the real 8-bit grit, supply a licensed 707 sample
set as an SFZ and load it through the device rack's ``Load SFZ'' entry;
\texttt{Bit Crusher} can add quantisation to an existing sound, but it cannot
turn a 909 into a 707.

\textbf{Boundary}: the 707's dryness is not ``turn the volume down''. It
comes from sample bit depth and envelope together. Muting a 909 open hat will
not produce a 707 closed hat.

\subsection{LinnDrum: the clean age of sampled drums}
\label{sec:iconic-linn}

Roger Linn built the LM-1 in 1980, the first drum machine to play back real
drum performances as digital samples. The 1982 LinnDrum (LM-2) brought the
price down and became a common sound in 1980s pop and R\&B. It used roughly
8-bit samples; the drums are clean and the hits are precise, a different path
from the 808's long analog tail. Michael Jackson's \emph{Thriller} album,
Prince's ``When Doves Cry'', and ``1999'' are repeatedly cited Linn-family
recordings.

\textbf{Ear fingerprint}: the kick is short and springy; the snare has a
clear hit without dirt; overall it sounds like a real kit recorded into a
machine and wiped clean. It has neither the 808's exaggerated low end nor the
707's obvious quantisation noise. It suits mid-tempo pop, R\&B, and sections
that need precise backbeats.

\textbf{Yeban landing}: approximate. Yeban has no Linn kit. The
\texttt{Acoustic R\&B} or \texttt{Acoustic Pop} style can supply a clean
backbeat and a mid-tempo frame, but that is not the 8-bit sample character of
1982. For a closer result, supply a licensed Linn sample SFZ, load it into the
sampler, and follow the ``short hit, dry, little low end'' direction in this
chapter. Do not force an 808 kit to do the job; its long tail gives the
substitution away immediately.

\textbf{Boundary}: the ``clean'' of a sample and the ``clean'' of a
synthesizer are not the same thing. Yeban's Acoustic style is another kit; it
must not be called a LinnDrum clone.

\subsection{MPC: sampler, slicer, and sequencer in one}
\label{sec:iconic-mpc}

Akai and Roger Linn introduced the MPC60 in 1988: 12-bit sampling, sixteen
velocity-sensitive pads, and a built-in sequencer in one box. It turned
``sample a drum break, slice it, rearrange it on the pads'' into a daily
workflow and became a centrepiece of 1990s hip-hop production. DJ Premier,
Pete Rock, and J Dilla are among the producers associated with the MPC line;
\emph{Ready to Die}, \emph{The Blueprint}, and \emph{Enter the Wu-Tang} are
common reference albums.

\textbf{Ear fingerprint}: it is not one drum sound but a feel: grainy
samples, slices cut on transients, and a performance with slight forward or
backward displacement. Pad velocity and sequencer micro-timing together make
the groove. Treating the MPC as only a sound pack throws away half of what it
is.

\begin{sloppypar}
\textbf{Yeban landing}: own material required; the slicing path exists.
Yeban has no sixteen-pad surface, but it has a batch slicing exit:
\texttt{yeban-app} \texttt{--slice-audio loop.wav} \texttt{--slice-out ./loop-slices} cuts a
drum loop at onsets into \texttt{<stem>\_<nnn>.wav} files and a
\texttt{<stem>.sfz} file. Load that SFZ into a sampler track through the
device rack, then rearrange it with MIDI notes. The rhythmic frame can use
\texttt{Acoustic Boom Bap} (default 90 BPM). For the MPC's slight drag, nudge
weak-beat notes in the piano roll or use the note's
\texttt{micro\_timing\_ticks} field; the current GUI has no dedicated MPC
swing panel.
\end{sloppypar}

\textbf{Boundary}: slicing is an audio-file operation, not pad performance
inside the project. Yeban has no MPC 16-pad surface and no ``MPC swing''
knob. Sample material must be licensed; do not slice a commercial recording
and publish it.

\section{Synthesizers}
\label{sec:iconic-synths}

Six instruments represent this line: the 303 for ``one line plus one filter''
minimalism; the Minimoog for the three-oscillator, ladder-filter analog
standard; the DX7 for digital FM; the Juno for affordable DCO polyphony; the
Prophet for programmable polyphony; and the Mellotron for the tape-sample
keyboard. Their technologies differ, and their sounds cannot replace one
another.

\subsection{TB-303: designed for bassists, rewritten by the dance floor}
\label{sec:iconic-303}

Roland introduced the TB-303 in 1981 as a portable ``computer bass'' for
guitarists and solo players: a saw or square wave into a 24 dB low-pass, with
a built-in step sequencer. It failed to convince its target users and was
discontinued in 1984 after about 10,000 units. In the mid-to-late 1980s,
Chicago producers used second-hand 303s to make the high-resonance ``acid''
sound; ``Acid Tracks'' (Phuture, 1987) is usually treated as the starting
point of acid house.

\textbf{Ear fingerprint}: monophonic, narrow, nasal saw or square; push the
cutoff and resonance to extremes and the note ``sings'', even self-oscillates;
slide makes the pitch travel instead of jump; overdrive or distortion makes
it more aggressive. All three together are the acid, not the word ``acid''
alone.

\textbf{Yeban landing}: approximate. Choose \texttt{Sub Bass}
(\texttt{synth:bass}); its low cutoff, higher resonance, and small drive give
the low-frequency frame of a 303. But Yeban's PolySynth has a fixed ADSR, no
clickable glide handle, and no dedicated 303 filter sweep. The project model
has a \texttt{slide} field, but the current GUI does not expose it as a
draggable glide control; true glide requires data/MCP writing or an external
303/VST3. Do not look for a ``303 glide'' knob in the GUI.

\textbf{Boundary}: Sub Bass is not a 303 clone. It lacks the 24 dB 303 filter
response and the accent's envelope retrigger. Lowering Sub Bass's cutoff gives
``dark'', not the scream of acid.

\subsection{Minimoog: portable, three oscillators, ladder low-pass}
\label{sec:iconic-minimoog}

Moog introduced the Minimoog in 1970, compressing the room-sized modular Moog
into a portable three-oscillator monophonic synth. Its ladder low-pass and
saturating mixer became the reference for what an analog synth should sound
like. About 13,000 were sold; Herbie Hancock, Stevie Wonder, Kraftwerk, and
Pink Floyd used it. ``Chameleon'''s bass, ``Autobahn'', and ``Shine On You
Crazy Diamond'' are common listening references.

\textbf{Ear fingerprint}: three oscillators can detune against one another
and stack into a fat bass or lead; the ladder low-pass resonance is not as
sharp as the 303's, but has a rounder corner and a little compression; push
the mixer and it saturates naturally. Its bass is not pure sub; it has
low-mid ``meat''.

\textbf{Yeban landing}: approximate. For bass, use \texttt{Sub Bass}
(\texttt{synth:bass}); for lead or mid-range melody, use
\texttt{Hollow Lead} (\texttt{synth:hollow}), whose two slightly detuned
branches simulate part of the thickness. Yeban has only two oscillator
branches, no third oscillator, no separate mixer saturation, and no
adjustable filter envelope. Do not write ``three oscillators'' as a Yeban
parameter.

\textbf{Boundary}: the Minimoog's fatness comes from three oscillators
interacting with a ladder filter; raising two wavetable levels will not
produce it. If you need a thicker three-layer sound, layer tracks, but watch
the phase relationship between low and mid frequencies.

\subsection{DX7: digital FM glass and bells}
\label{sec:iconic-dx7}

Yamaha introduced the DX7 in 1983. It used digital FM synthesis with six
operators and thirty-two algorithms, and was the first digital synthesizer to
reach recording studios in large numbers. It had no traditional knobs;
programming was menu-driven, but its preset sounds became a common language
of 1980s pop. The Koto in ``When Doves Cry'', the Bass 1 in ``Take On Me'',
``Smooth Operator'', and ``Smooth Criminal'' are commonly cited DX7 examples.

\textbf{Ear fingerprint}: glass, metal, bells, and the ``ding'' of an
electric piano; the sound can be clean or full of complex inharmonic content.
Its envelope and algorithm combinations produce a ``digital cold light'' that
analog synths cannot. The DX7's E Piano 1 preset even displaced part of the
Rhodes market.

\textbf{Yeban landing}: approximate, and distinguish the name from the
engine. Yeban's \texttt{FM Bell} (\texttt{synth:metallic}) has FM in its name,
but the built-in engine is a two-oscillator wavetable into a ladder low-pass,
not an FM operator engine. It can give a glass/metal impression; it cannot
give DX7 algorithms, operator ratios, or feedback. If
\texttt{TX81Z - FM Piano.sfz} is installed, load it into a sampler track for a
four-operator FM electric-piano approximation; the TX81Z is not a DX7, but it
belongs to the same FM family. For bells use \texttt{FM Bell}; for short
glass plucks use \texttt{Glass Pluck} (\texttt{synth:glass}).

\textbf{Boundary}: do not treat the \texttt{FM Bell} parameter table as a DX7
algorithm table. Yeban has no operators, ratios, feedback, or algorithm
selection. ``FM Bell'' is only the name of a wavetable preset.

\subsection{Juno: DCOs, a sub oscillator, and ``that chorus''}
\label{sec:iconic-juno}

Roland introduced the Juno-60 in 1982 and the Juno-106 in 1984. They used
digitally controlled oscillators (DCOs) for stable tuning, a sub oscillator
to thicken the bass, and a built-in chorus to widen the single-oscillator
sound. The Juno made polyphonic synths affordable and became a common voice
in 1980s synth-pop and house; ``Take On Me'', ``Time After Time'', and ``Can
You Feel It?'' used the Juno family.

\textbf{Ear fingerprint}: DCO cleanliness and stability, plus the sub
oscillator's low end, plus the width and softness of the chorus. Pads are
``wide without mush''; plucks are ``bright without hardness''. The chorus is
not decoration; it is part of the timbre's skeleton.

\textbf{Yeban landing}: approximate, with the chorus missing. Use
\texttt{Choir Pad} (\texttt{synth:vocal}) for a wide pad and
\texttt{Glass Pluck} (\texttt{synth:glass}) for short plucks. Yeban has no
built-in chorus, flanger, or phaser, so the Juno width can only be assembled
from layering, detune, reverb, and tape saturation. That is not a Juno clone.
For a real chorus, use a third-party VST3 or external production.

\textbf{Boundary}: the Juno sub oscillator is not an adjustable Yeban
parameter; Choir Pad's vocal wavetable is not a sampled chorus. The two can
correspond in role, not in circuit.

\subsection{Prophet-5: the first programmable polyphonic synth}
\label{sec:iconic-prophet}

Sequential Circuits introduced the Prophet-5 in 1978. It used a
microprocessor to store patches and was the first fully programmable
polyphonic synthesizer: five voices, two oscillators, a ladder/Curtis low-
pass, and Poly-Mod. It became a staple of 1980s pop and film music;
``In the Air Tonight'', the title track of \emph{Thriller}, ``Everything In
Its Right Place'', and ``I Can't Go for That'' are common references.

\textbf{Ear fingerprint}: five-voice polyphony, warm oscillators with bite,
and storable presets. Pads have a choral quality without relying on an effect;
leads have analog thickness; the bass can also sink. Poly-Mod lets the filter
and oscillators modulate each other for atypical sweeps.

\textbf{Yeban landing}: approximate. Use \texttt{Hollow Lead} for a lead,
\texttt{Choir Pad} for a pad, and \texttt{Warm Organ} for sustained chords.
Yeban's sixteen voices exceed the Prophet's five, but it has no Poly-Mod,
no independent filter envelope, and no oscillator sync. Do not call it a
Prophet re-creation.

\textbf{Boundary}: more polyphony does not make it sound like a Prophet. The
Prophet's warmth comes from analog oscillators and filters; Yeban's
wavetables can supply only role and approximate spectrum.

\subsection{Mellotron: a ``sampling keyboard'' made of tape}
\label{sec:iconic-mellotron}

The Mellotron appeared in Birmingham, England, in 1963. Each key triggered a
magnetic tape that played a pre-recorded instrument; each tape lasted about
eight seconds, played while held, and rewound on release. It was not a modern
digital sampler but a mechanical tape-sample keyboard. In the 1960s the
Beatles used it on ``Strawberry Fields Forever'', King Crimson on \emph{In
the Court of the Crimson King}, and it later became a progressive-rock
machine for strings, choir, and flute.

\textbf{Ear fingerprint}: strings, choir, and flute carry tape noise, slight
wow, and the unevenness of ``one tape per key''. It is not clean sampling;
it is recorded playback with mechanical traces. The eight-second limit also
means it cannot sustain forever.

\textbf{Yeban landing}: own material or approximation. Yeban has no
Mellotron tape samples. \texttt{Choir Pad} can make a choral bed, and
\texttt{Tape Saturation}'s \texttt{tape\_wow}/\texttt{tape\_flutter} can add
transport wobble to the whole signal, but that is not per-key Mellotron tape.
For a closer result, supply a licensed Mellotron sample SFZ. Do not call
Choir Pad a Mellotron.

\textbf{Boundary}: Tape Saturation's wow/flutter modulates the entire signal,
including the dry sound; it is not per-key Mellotron tape and cannot produce
an eight-second tape break.

\section{Electric instruments}
\label{sec:iconic-electric}

Electric instruments are not ``a preset on a synthesizer''. Their sound
mechanisms differ: Rhodes uses metal tines, Wurlitzer uses metal reeds, and
Hammond uses tone wheels and drawbars. The mechanism decides the transient
and the harmonics; a wavetable name cannot replace it.

\subsection{Rhodes: tines, bells, and stereo tremolo}
\label{sec:iconic-rhodes}

Harold Rhodes's design entered production as the Fender Rhodes in 1965. A
hammer strikes a metal tine and a pickup converts the vibration into an
electrical signal. The sound is bell-like and electric-piano-like, with a
long, round tail; Suitcase models added stereo tremolo. Miles Davis's
\emph{In a Silent Way} and \emph{Bitches Brew}, the Doors' ``Riders on the
Storm'', and Stevie Wonder's ``You Are the Sunshine of My Life'' are classic
references.

\textbf{Ear fingerprint}: a bell-like ``ding'' at the attack, metallic
overtones in the mid-highs, and a tail longer than a Wurlitzer but more
``electric'' than an acoustic piano. Stereo tremolo makes the left and right
channels swing slightly out of phase. It suits mid-tempo ballads, fusion, and
keyboard layers that need ``warm but electric''.

\textbf{Yeban landing}: approximate. \texttt{FM Bell}
(\texttt{synth:metallic}) can give a bell-like attack and metallic tail, but
its fixed ADSR and wavetable are not a tine. If
\texttt{TX81Z - FM Piano.sfz} is installed, loading it into a sampler track
gets closer to an FM electric piano. Yeban has no tremolo/stereo-tremolo
device and no Rhodes-specific sample; for a real Rhodes, supply a licensed
SFZ or a third-party instrument.

\textbf{Boundary}: lowering FM Bell's low end does not make it a Rhodes. The
tine's inharmonic content and pickup response cannot be fully reproduced by a
wavetable.

\subsection{Wurlitzer: reeds, bite, and bark}
\label{sec:iconic-wurlitzer}

The Wurlitzer electric piano was produced from 1954 to 1983. It used metal
reeds and sounded harder and more direct than the Rhodes; the 200A is the most
common stage model. Ray Charles's ``What'd I Say'', Joe Zawinul's ``Mercy,
Mercy, Mercy'', Supertramp's ``The Logical Song'', and Queen's ``You're My
Best Friend'' are common references.

\textbf{Ear fingerprint}: a more biting attack than the Rhodes, stronger
midrange, and a shorter tail. It has a reed-like ``bark'', especially in the
low register. It suits soul, R\&B, rock, and sections that need an electric
piano with no float.

\textbf{Yeban landing}: approximate or own material. Yeban has no Wurlitzer
reed model. \texttt{FM Bell} with a lower cutoff and a little
\texttt{Tape Saturation} drive can approach the bite, but there is no reed
mechanism and no Tremolo. To keep the 200A's bite, supply a licensed
Wurlitzer sample SFZ.

\textbf{Boundary}: both are electric pianos, but one is a tine and the other
a reed. Brightening a Rhodes bell does not produce a Wurlitzer bark.

\subsection{Hammond: tone wheels, drawbars, and Leslie}
\label{sec:iconic-hammond}

The Hammond B-3 was produced from 1955 to 1974. Tone wheels generate
near-sine harmonics, and nine drawbars mix them in proportion. The Leslie
rotary speaker uses a two-speed rotating horn to create Doppler and amplitude
modulation. ``Green Onions'' (Booker T. \& the M.G.'s, Hammond M-3), ``A
Whiter Shade of Pale'' (Procol Harum), ``Like a Rolling Stone'' (Bob Dylan),
and ``I Want You (She's So Heavy)'' (The Beatles) are Hammond-family
references.

\textbf{Ear fingerprint}: tone wheels give a stable harmonic stack; drawbars
decide which overtone is louder; the Leslie makes the sound rotate in space,
slow like breathing and fast like a sweep. It has no piano attack transient;
press a key and the organ tone sustains.

\begin{sloppypar}
\textbf{Yeban landing}: approximate, with drawbars and Leslie missing. Use
\texttt{Warm Organ} (\texttt{synth:organ}) for a sustained organ timbre. If
\texttt{DrawbarOrganEmulation}\allowbreak\texttt{-20190712.sfz} is installed, load it into a
sampler track for a drawbar-organ sample that is closer. Yeban has no
drawbar controls and no rotary-speaker device; for a Leslie, use a
third-party VST3 or external production.
\end{sloppypar}

\textbf{Boundary}: Warm Organ is a fixed-wavetable organ preset, not a drawbar
organ. It has no nine-drawbar ratio control and no Leslie rotating filter.

\section{Six ``whys''}
\label{sec:iconic-why}

This section separates six principles that are easy to misuse in production.
They are not historical trivia; they are the mechanisms that decide why a
sound can be used a certain way.

\subsection{Why an 808 kick can be a bass}
\label{sec:iconic-why-808}

The 808 kick is not a short ``thump''; it is a pulse with a falling pitch
entering a long decay. Its fundamental usually sits in the tens of hertz, and
the decay is long enough that repetition forms a continuous low-frequency
line. That line is the bassline. A drum-machine kick is a rhythmic event, but
when its tail crosses the next beat, the event becomes a sustained bass.

To use this well in Yeban, separate the low-frequency carrier: either let the
kick's long tail be the bass and keep \texttt{bass808} as decoration, or let
\texttt{bass808} be the sustained line and shorten the kick to an impact.
With both fully open, the low end masks itself; it sounds muddier, not thicker.
Check by soloing the low tracks and listening whether the kick tail and
\texttt{bass808} fight for the same peak.

\subsection{Why a 303 glide sounds acid}
\label{sec:iconic-why-303}

A 303 glide is not decoration: it makes the pitch travel continuously between
notes. At the same time, a high-resonance low-pass pushes certain harmonics
forward, so the pitch change becomes ``the filter singing''. With glide,
resonance, and overdrive together, the timbre acquires an unstable, liquid
``acid'' quality; glide alone or resonance alone is not enough.

Yeban has no 303 glide handle. The project model has a \texttt{slide} field,
but the current GUI does not expose a draggable glide control. Sub Bass can
supply a low-frequency frame, not the 303 filter sweep. To study the principle,
first listen to a reference for two things: the pitch sliding and the resonant
peak moving with it. Then decide whether an external 303/VST3 is worth it. Do
not force Sub Bass to imitate it.

\subsection{Why slicing and swing can destroy each other}
\label{sec:iconic-why-slicing}

Slicing cuts an audio file at transients; swing pushes even beats or weak
beats later. Both change time, but at different layers: slicing changes the
boundaries inside the audio file, swing changes note trigger times. Slice
first and swing later: the slices themselves do not move, only the triggers
do. If a slice boundary falls in the middle of a drum tail, the new trigger
after swing cuts the tail and the result sounds more chopped, not more
groovy.

Yeban's slicing exit is a batch process: \texttt{--slice-audio} generates
one WAV per onset and an SFZ, and the sampler plays them from MIDI notes.
Swing or micro-timing belongs on the MIDI trigger layer, not in an expectation
that slices will rearrange themselves. To preserve a drum loop's groove, put
the cuts between transients; to rearrange the groove, cut cleanly first, then
move triggers in the piano roll.

\subsection{Why 707 and 909 metal sound different}
\label{sec:iconic-why-707-909}

The 707's metal is all sampled, mostly 8-bit, so it is dry, short, and
slightly noisy from quantisation. The 909's metal is 6-bit sampled, while its
kick and toms remain analog, so the metal is cleaner and more clipped while
the kick is harder. The difference is not ``which is more advanced''; it is
the sample bit depth and the ratio of analog to digital.

Yeban has no 707 kit. Using the 909 kit's short closed hats to approximate the
707's dryness covers only part of it. To practise the distinction, take the
same four-on-the-floor pattern, mute the 909 open hat and shorten the closed
hat, then compare with a 707 reference: the 707 closed hat is shorter and more
``square'', while the 909 metal has more cymbal tail.

\subsection{Why Rhodes and Wurlitzer are both electric pianos but unlike}
\label{sec:iconic-why-rhodes-wurli}

Rhodes uses a metal tine, so the attack is bell-like, the tail is long, and
the overtones lean metallic. Wurlitzer uses a metal reed, so the attack bites,
the tail is shorter, and the midrange is stronger. Different mechanisms mean
different transients and harmonics; this is not something ``brighten it'' or
``add a little distortion'' can swap.

Yeban has no dedicated samples for either. FM Bell can give the Rhodes bell
attack, but not the tine's inharmonic content; brightening FM Bell will not
produce the Wurlitzer reed bite. If you need one of them, prepare a licensed
sample of that instrument; do not use one electric-piano preset to impersonate
the other.

\subsection{Why Hammond drawbars and Leslie must be heard separately}
\label{sec:iconic-why-hammond}

Drawbars decide ``which overtones are louder'', a frequency-domain mix;
the Leslie decides ``how the sound rotates in space'', a time/amplitude
modulation. The two are independent: a Hammond without a Leslie is still a
Hammond; a ``rotating organ'' without drawbars is just a rotating pipe organ.
Yeban's Warm Organ gives only a fixed organ wavetable, with no drawbar ratios
and no Leslie. When listening to references, first hear the drawbar harmonic
mix, then the Leslie's slow/fast rotation separately; do not collapse them
into one ``organ flavour''.

\section{Yeban mapping table}
\label{sec:iconic-yeban-map}

The table compresses the fourteen instruments into one lookup. ``Directly
usable'' means a factory style or preset carries the main traits;
``approximate'' means only role and part of the spectrum can be supplied;
``own material'' means Yeban does not ship a licensed sample and the user
must prepare an SFZ or a third-party instrument.

\begin{longtable}{p{0.13\linewidth}p{0.07\linewidth}p{0.20\linewidth}p{0.31\linewidth}p{0.15\linewidth}}
\hline
Instrument & Year & Ear fingerprint & Yeban landing & Grade \\
\hline
\endhead
TR-808 & 1980 & long-tail kick, noisy snare, analog percussion &
\texttt{Trap}/\texttt{Hip-Hop} style + \texttt{bass808} & direct \\
TR-909 & 1983 & hard kick, cracked snare, sampled metal &
\texttt{House}/\texttt{Techno}/\texttt{EDM} style & approximate \\
TR-707 & 1985 & dry, short, 8-bit metal & 909 short closed hats; or own
707 SFZ & own material \\
LinnDrum & 1982 & clean samples, short hits, no long tail &
\texttt{Acoustic R\&B}/\texttt{Acoustic Pop}; or own SFZ & own material \\
MPC & 1988 & grainy samples, slicing, micro-timing &
\texttt{--slice-audio} $\to$ SFZ $\to$ sampler + \texttt{Boom Bap} & own
material \\
TB-303 & 1981 & mono, high resonance, glide & \texttt{Sub Bass}; no glide/GUI
sweep & approximate \\
Minimoog & 1970 & three oscillators, fat bass, ladder low-pass &
\texttt{Sub Bass} / \texttt{Hollow Lead} & approximate \\
DX7 & 1983 & glass, metal, FM electric piano & \texttt{FM Bell} /
\texttt{Glass Pluck}; or \texttt{TX81Z} SFZ & approximate \\
Juno & 1982 & DCO, sub oscillator, chorus & \texttt{Choir Pad} /
\texttt{Glass Pluck}; no chorus & approximate \\
Prophet-5 & 1978 & five voices, warm, Poly-Mod & \texttt{Hollow Lead} /
\texttt{Choir Pad} / \texttt{Warm Organ} & approximate \\
Mellotron & 1963 & tape strings/choir, transport traces & \texttt{Choir Pad} +
\texttt{Tape Saturation}; or SFZ & own material \\
Rhodes & 1965 & tine, bell, long tail & \texttt{FM Bell}; or \texttt{TX81Z}
SFZ & approximate \\
Wurlitzer & 1954 & reed, bite, short tail & \texttt{FM Bell} low cutoff +
drive; or SFZ & approximate \\
Hammond & 1955 & drawbar harmonics, Leslie rotation & \texttt{Warm Organ};
or \texttt{DrawbarOrganEmulation} SFZ & approximate \\
\hline
\end{longtable}

\section{Listening exercises}
\label{sec:iconic-exercises}

The exercises use only drum/keyboard passages from reference records; do not
sing along and do not copy lyrics. For each, listen to the reference first,
then make an A/B in Yeban. If you cannot hear the difference, go back to the
four dimensions; do not hide it by raising the volume.

\begin{enumerate}
  \item \textbf{808 versus 909 kick}: find an 808 long tail and a 909 short
        hit, two bars each. In Yeban, generate \texttt{Hip-Hop} and
        \texttt{House} styles, keep only the kick, and compare whether the
        tail crosses the beat.
  \item \textbf{303 glide}: choose a passage from ``Acid Tracks'' where the
        pitch travels continuously. Play the same pitches on Sub Bass without
        glide, then add \texttt{slide} (if using external data); hear the
        difference between ``jump'' and ``travel''.
  \item \textbf{707 versus 909 metal}: mute the 909 open hat, shorten the
        closed hat, and compare with a 707 reference. Note that the 707 closed
        hat is shorter and drier, while the 909 cymbal tail is longer.
  \item \textbf{Rhodes versus Wurlitzer}: in the same slow song, listen to the
        Rhodes in ``Riders on the Storm'' and the Wurlitzer in ``The Logical
        Song''. In Yeban use FM Bell for both and change only cutoff and
        drive; feel that ``bell'' and ``reed bite'' are not the same direction.
  \item \textbf{Hammond drawbars versus Leslie}: listen to the organ
        harmonics in ``Green Onions'', then a Hammond passage with Leslie
        rotation. In Yeban, use Warm Organ for sustained chords and admit that
        drawbars and rotation are missing; do not use reverb as a Leslie
        substitute.
\end{enumerate}

\section{Help-topic mapping for this chapter}
\label{sec:iconic-help-map}

\begin{longtable}{p{0.42\linewidth}p{0.46\linewidth}}
\hline
Section & In-app help anchor \\
\hline
\endhead
Four dimensions for listening & \texttt{iconic-listening} \\
TR-808 & \texttt{iconic-drum-808} \\
TR-909 & \texttt{iconic-drum-909} \\
TR-707 & \texttt{iconic-drum-707} \\
LinnDrum & \texttt{iconic-drum-linn} \\
MPC & \texttt{iconic-drum-mpc} \\
TB-303 & \texttt{iconic-synth-303} \\
Minimoog & \texttt{iconic-synth-minimoog} \\
DX7 & \texttt{iconic-synth-dx7} \\
Juno & \texttt{iconic-synth-juno} \\
Prophet-5 & \texttt{iconic-synth-prophet} \\
Mellotron & \texttt{iconic-synth-mellotron} \\
Rhodes & \texttt{iconic-electric-rhodes} \\
Wurlitzer & \texttt{iconic-electric-wurlitzer} \\
Hammond & \texttt{iconic-electric-hammond} \\
Six ``whys'' & \texttt{iconic-why} \\
Yeban mapping table & \texttt{iconic-yeban-map} \\
Listening exercises & \texttt{iconic-exercises} \\
\hline
\end{longtable}

\section{Source boundaries and further reading}
\label{sec:iconic-sources}

Maker, year, production figures, and work attribution in this chapter prefer
official manufacturer history pages, product pages, and verifiable
encyclopaedic entries. Tone descriptions are editorial judgments, not
measurements. Further reading can check Roland's official 808/909/303 history
pages, Yamaha's DX7 history page, Sequential's Prophet-5 page, Streetly
Electronics' Mellotron material, and Fender Rhodes, Wurlitzer, and Hammond
product-history material. Works are named by title only; lyrics are not quoted.
If lyrics are needed, handle the copyright separately and do not put lyrics
into this manual.

\begin{itemize}[leftmargin=1.5em]
\item \url{https://www.roland.com/global/} — Roland official site: check the manufacturer history and product material for the 808, 909, 303, and other iconic instruments.
\item \url{https://www.sequential.com/} — Sequential official site: check the Prophet model lineage and design background.
\end{itemize}
